% Solution: Tadelis (2013) Exercise 16.5 "Drug Ads". Input by \solitem. Check: check_T1.py.
\begin{exsol}{Tadelis Exercise 16.5}{337}{Foundations}{3}
\asked (a) The extensive form; (b) a separating PBE; (c) a pooling PBE; (d) what changes when
$R > r > c > 0$ instead of $R > c > r > 0$.

\method The five-step method: guess a profile, Bayes' rule, the consumer's reply, the company's
incentive check, off-path beliefs (Tadelis \S16.1, pp.~319--323). The consumer buys iff his belief
that the drug is effective is at least $\frac12$: buying pays $\mu - (1 - \mu) = 2\mu - 1$ against $0$.

\solution \textbf{(a) Tree.} Nature picks $H$ or $L$ with probability $\frac12$; the company chooses
$A$ (cost $c$) or $N$; the consumer, seeing only the message, buys ($B$) or not ($\bar B$).
Payoffs (company, consumer): after $H$: $(R - c, 1)$, $(-c, 0)$ for $A$ and $(R, 1)$, $(0, 0)$ for
$N$; after $L$: $(r - c, -1)$, $(-c, 0)$ for $A$ and $(r, -1)$, $(0, 0)$ for $N$.
\sigtreenames{$H$}{$L$}{$A$}{$N$}{$B$}{$\bar B$}{$\frac12$ each}
\sigtree{R-c,1}{-c,0}{r-c,-1}{-c,0}{R,1}{0,0}{r,-1}{0,0}
(The left branches are message $A$, the right branches message $N$; upper node = type $H$, replies
$B$ above and $\bar B$ below.)

\textbf{(b) Separating PBE ($R > c > r$).} Strategies: $H$ plays $A$, $L$ plays $N$. Beliefs:
$\mu(A) = 1$, $\mu(N) = 0$. Replies: buy after $A$ ($\mu = 1 \ge \frac12$), do not buy after $N$.
Incentives: $H$ gets $R - c > 0$ against $0$ from $N$ (since $R > c$); $L$ gets $0$ from $N$ against
$r - c < 0$ from imitating $A$ (since $c > r$). There are no off-path messages. The reverse
profile ($H$: $N$, $L$: $A$) fails: $\mu(N) = 1$ makes the consumer buy, so $L$ would gain $r > -c$
by switching to $N$.

\textbf{(c) Pooling PBE ($R > c > r$).} Pooling on $A$ is not an equilibrium: $L$ earns at most
$r - c < 0$ (the consumer buys at best, since $\mu(A) = \frac12$), while deviating to $N$ earns him
at least $0$ whatever the consumer does; no off-path belief can prevent this.
Pooling on $N$ is: both types play $N$, $\mu(N) = \frac12$, the consumer is indifferent and buys
(or randomizes). The types get $R$ and $r$. A deviation to $A$ gives $H$ at most $R - c < R$ and $L$
at most $r - c < r$, whatever the consumer does after $A$, so every off-path belief $\mu(A)$ works.
PBE: both $N$; consumer buys after $N$; after $A$ he best-replies to any belief $\mu(A) \in [0,1]$.
(If he does not buy after $N$ the types earn $0$; the profile is still an equilibrium provided he does
not buy after $A$ either, i.e.\ $\mu(A) \le \frac12$.)

\textbf{(d) $R > r > c > 0$.} Advertising is now cheaper than the low type's revenue from a sale,
so $L$ would mimic: in $(A_H, N_L)$ the consumer buys after $A$, and $L$ earns $r - c > 0$ from
imitating against $0$. There is no separating PBE (the reverse profile fails as before). Pooling on
$N$ still works, and now pooling on $A$ works too: $\mu(A) = \frac12$, the consumer buys, the types
get $R - c$ and $r - c > 0$, and a deviation to $N$ is deterred by the off-path belief
$\mu(N) < \frac12$ (no purchase, payoff $0 < r - c$). Advertising stops separating because it is
cheap relative to the low type's gain: single crossing fails.

\verify \texttt{check\_T1.py} enumerates all pure profiles for $(R, c, r) = (10, 6, 3)$ and
$(10, 2, 3)$ on a grid of beliefs of step $\frac1{60}$: the first has the separating PBE $(A_H, N_L)$ and
pooling only on $N$; the second has no separating PBE and pooling on $A$ and on $N$. Tadelis prints
no answer.

\answer{(b) $H$ advertises, $L$ does not; $\mu(A) = 1$, $\mu(N) = 0$; buy after $A$ only.
(c) Both $N$, buy after $N$, any belief after $A$ (pooling on $A$ is not an equilibrium).
(d) No separating PBE; pooling on $N$ and also on $A$ (with $\mu(N) < \frac12$ after the off-path $N$).}
\end{exsol}
